\scanpage{68}\mspage{33}
$\varepsilon,\varepsilon',\varepsilon''\in\{\pm1\}$ tels que l'on ait $\varepsilon x+\varepsilon'y+\varepsilon''z=0$ ou encore qu'on a une égalité de la forme
\[
 z=\deleted{\unsure{\pm}}\varepsilon x+\varepsilon'y\qquad(\varepsilon,\varepsilon'\in\{\pm1_\Lambda\})\tag*{\deleted{(179)} $(*)$}
\]
ce qui est évidemment faux. Par contre, \deleted{\unsure{pour tout $i\in$|pour tout si $i\in$}} \inserted{si on se donne}, en plus de $A_M$, l'application $\tau_A:A_M\longrightarrow\unsure{\pi|\Pi}_M$, i.e la décomposition de $A_M$ suivant les types $i\in\pi$, on réconstitue $F_M$ comme
\[
 \left\{\begin{aligned}
 F_M&=\coprod_{i\in\pi}F_{M,i}\\
 F_{M,i}&=\{f\in\mathfrak P_3(S_M)\mid\forall a\in\mathfrak P_2(f),\text{ on a }a\in A_{M,i}\}
 \end{aligned}\right.\tag*{(\unsure{178|180})}
\]
cette dernière relation en effet signifie que si $x,y,z\in M$ sont tels que
\[
 x\wedge y,\ y\wedge z,\ z\wedge x\quad\in i=\{e,-e\}\subset(\det M)^*
\]
alors on doit avoir une relation $(*)$ ci dessus, en effet. \unsure{comme|Si} $(x,y)\in\operatorname{Rep}(M)$, on aura $z=\varepsilon x+\varepsilon'y$, \inserted{\unsure{avec}} $\varepsilon,\varepsilon'\in\Lambda$, et on aura, en multipliant cette relation par $x$
\[
 \varepsilon'x\wedge y=x\wedge z\qquad\text{d'où}\quad \varepsilon'\in\{\pm1_\Lambda\}
\]
et de \unsure{$=$|même} $\varepsilon\in\{\pm1_\Lambda\}$, cqfd.

Dans le cas \deleted{\unsure{B}}A) par contre $(2.1_\Lambda=0)$,
% archive sheet number: 67
